\chapter{Introduction}
\label{chap:introduction}

\section{Motivation}

Large language models (LLMs) are increasingly used as components of
multi-agent systems in which models critique, debate, verify, or revise one
another's outputs. Such systems can improve factuality and reasoning by
combining multiple model responses~\cite{du2024debate}. Their reliability,
however, depends on more than whether a Critic can generate a plausible
recommendation. The receiving model must determine whether that recommendation
should change its answer. Blind adoption can propagate incorrect feedback,
whereas indiscriminate resistance can prevent useful correction. Research on
self-correction and sycophancy has shown that a model's willingness to revise
is not equivalent to reliable error correction
~\cite{huang2024selfcorrect,sharma2024sycophancy}.

Behavioral evaluation can determine whether a Solver changed its answer, but
it cannot explain how the Solver internally represented the feedback or which
parts of that representation contributed to its decision. A hidden activation
may correlate with Adoption because it represents disagreement, uncertainty,
answer identity, presentation order, or merely a downstream consequence of
the decision. This creates a central problem for mechanistic interpretability:
finding a predictive internal signal is not the same as identifying a feature
with a coherent meaning or demonstrating that the signal participates in the
model's computation.

Sparse autoencoders (SAEs) provide one approach to this problem. An SAE
decomposes a dense model activation into a larger collection of sparsely active
features. Compared with individual residual-stream dimensions, these features
may isolate more coherent patterns in the model's representation
~\cite{cunningham2023sparse,gao2025scaling}. Sparsity alone, however, does not
guarantee interpretability. A useful feature interpretation should be
supported by semantically coherent activating examples, generalize to held-out
interactions, distinguish the target behavior from nuisance variables, and
make correct predictions about what will happen when the feature is modified.

This project studies feature interpretability through the concrete behavioral
case of feedback uptake. A Gemma Solver first answers a multi-hop question, a
Qwen Critic independently evaluates the problem, and the Solver then chooses
between its original answer and the Critic-supported answer. This controlled
candidate-presentation design makes Adoption and Retention directly observable
while preserving Agreement and unresolved outputs as separate outcomes. We
then examine the Solver's residual-stream activations with SAEs, predictive
probes, feature-level analyses, and causal interventions. Feedback uptake is
therefore not only the behavior being modeled; it is a testbed for asking
whether sparse features provide explanations that are semantically meaningful,
behaviorally relevant, and causally faithful.

\section{Research Question and Objectives}

The primary research question is:

\begin{quote}
Can sparse autoencoders identify human-interpretable features in a language
model's residual stream that explain and causally influence whether the model
adopts or resists feedback?
\end{quote}

We address this question through four objectives:

\begin{enumerate}
    \item Build a controlled Solver--Critic interaction pipeline and a
    leakage-aware dataset of behavioral outcomes and Solver activations.

    \item Represent feedback-processing activations using pretrained and
    interaction-specific SAEs, and evaluate whether those representations
    preserve the underlying residual-stream information with useful sparsity.

    \item Identify and interpret sparse features that predict Adoption versus
    Retention on held-out data, using activating examples, counterexamples,
    control probes, and random-feature baselines to evaluate their meaning and
    specificity.

    \item Test whether selected sparse features are causally involved in
    feedback uptake by suppressing or amplifying them during held-out forward
    passes and comparing the resulting
    behavioral changes with reconstruction, random-feature, and equal-size
    activation controls.
\end{enumerate}

These objectives separate several claims that are often conflated. SAE
reconstruction evaluates whether the representation preserves information;
probing evaluates whether that information predicts a behavior; semantic
analysis evaluates whether a feature admits a stable human-readable
description; and intervention evaluates whether the feature or direction can
influence the behavior. Evidence at one level does not automatically establish
the others.

\section{Contributions}

This project makes the following contributions:

\begin{itemize}
    \item A reproducible Solver--Critic pipeline that records feedback-uptake
    behavior and residual-stream activations under a controlled
    candidate-presentation protocol.

    \item A labeled interaction dataset that distinguishes Adoption,
    Retention, Agreement, and outputs that cannot support a reliable binary
    interpretation, with question-family grouping to reduce data leakage.

    \item A comparison of sparse representations derived from pretrained,
    custom-trained, and adapted SAEs at selected Solver layers.

    \item A feature-discovery workflow combining regularized probes,
    selectivity checks, random-feature controls, and examination of dominant
    features rather than treating probe coefficients as explanations by
    themselves.

    \item Held-out causal evaluations that compare targeted suppression and
    amplification with norm-matched random edits, SAE reconstruction controls,
    and raw-activation baselines while separating intended behavioral switches
    from broken or unclassifiable generations.

    \item An explicit account of the limits of feature interpretation: causal
    effectiveness supports behavioral relevance, but human interpretability
    additionally requires coherent semantic evidence and robustness across
    examples and controls.
\end{itemize}

\section{Report Roadmap}

Chapter~2 introduces transformer residual streams, mechanistic
interpretability, sparse autoencoders, and LLM feedback behavior. Chapter~3
reviews related work on SAE feature evaluation, representation engineering,
sycophancy, and multi-hop question answering. Chapters~4 and~5 describe the
models, MuSiQue data, exploratory analysis, synthetic-data development, and
the final interaction dataset. Chapter~6 presents the experimental pipeline,
labeling procedure, activation collection, SAE training, feature-selection
protocol, interpretation criteria, causal interventions, and statistical
controls. Chapter~7 reports the behavioral, representational, predictive, and
causal results. Chapter~8 evaluates what those results do and do not establish
about feature interpretability, discusses threats to validity and limitations,
and identifies directions for future work. Chapter~9 concludes by answering
the research question in light of the combined evidence.
